% \documentclass[12pt]{article}
% \setlength{\textwidth}{17cm}
% \setlength{\textheight}{24cm}
% \setlength{\topmargin}{-2cm}
% \setlength{\footskip}{1cm}
% \setlength{\evensidemargin}{0cm}
% \setlength{\oddsidemargin}{0cm}
% \setlength{\parindent}{0cm}

% \usepackage{allrunes}
% \usepackage{amsmath}
% \usepackage[magyar]{babel}
% \usepackage[T1]{fontenc}
% \usepackage[utf8]{inputenc}
% \usepackage{fixltx2e}
% \usepackage{multirow}
% \usepackage{hyperref}
% \usepackage{amsfonts}
% \usepackage{amsthm}
% \usepackage{braket}

% \theoremstyle{plain}
% \usepackage{graphicx}

% %\usepackage{gensymb}
% \usepackage{float}

% %linkek 
% \newcommand{\ndurl}[2]{\textit{\href{#1}{\underline{#2}}}}

% %% New commands
% %\newcommand{\ket}[1]{\left| #1 \right >}
% %\newcommand{\bra}[1]{\left < #1 \right |}
% %\newcommand{\braket}[2]{\left < #1 \middle | #2 \right>}
% \newcommand{\norm}[1]{\left{||} #1 \right{||}}
% \newcommand{\commut}[2]{\left [ #1 , #2 \right]}

% %% Pauli matrices
% \newcommand{\sigx}{\sigma_x}
% \newcommand{\sigy}{\sigma_y}
% \newcommand{\sigz}{\sigma_z}

% \newcommand{\paulix}{
%     \left( \begin{array}{cc}
%         0 & 1 \\
%         1 & 0
%     \end{array}
%     \right)
% }
% \newcommand{\pauliy}{
%     \left( \begin{array}{cc}
%         0 & -i \\
%         i & 0
%     \end{array}
%     \right)
% }
% \newcommand{\pauliz}{
%     \left( \begin{array}{cc}
%         1 & 0 \\
%         0 & -1
%     \end{array}
%     \right)
% }

% \newcommand{\dd}{\mathrm{d}}

% \newtheorem*{theorem*}{Tétel}
% \newtheorem*{def*}{Definíció}
% \newtheorem*{pld*}{Példa}
% \newtheorem*{megj}{Megjegyzés}


% \begin{document}
% \title{Kvantummechanika gyorstalpaló}
% \maketitle
% \newpage
\documentclass[12pt]{article}
\usepackage[utf8]{inputenc}

\setlength{\textwidth}{17cm}
\setlength{\textheight}{24cm}
\setlength{\topmargin}{-2cm}
\setlength{\footskip}{1cm}
\setlength{\evensidemargin}{0cm}
\setlength{\oddsidemargin}{0cm}
\setlength{\parskip}{1em}
\setlength{\parindent}{2em}

\usepackage{allrunes}
\usepackage{amsmath}
\usepackage[magyar]{babel}
\usepackage[T1]{fontenc}
\usepackage{fixltx2e}
\usepackage{multirow}
\usepackage{hyperref}
\usepackage{amsfonts}
\usepackage{amsthm}
\usepackage{amssymb}

\theoremstyle{plain}
\usepackage{graphicx}

\newcommand{\code}[1]{\texttt{#1}}

%\usepackage{gensymb}
\usepackage{float}

%linkek 
\newcommand{\ndurl}[2]{\textit{\href{#1}{\underline{#2}}}}

\usepackage{braket}

\newcommand{\dd}{\textrm{d}}

\newcommand{\mev}{~\textrm{MeV}}
\newcommand{\fm}{~\textrm{fm}}

%% Pauli matrices
\newcommand{\sigx}{\sigma_x}
\newcommand{\sigy}{\sigma_y}
\newcommand{\sigz}{\sigma_z}

\newcommand{\paulix}{
    \left( \begin{array}{cc}
        0 & 1 \\
        1 & 0
    \end{array}
    \right)
}
\newcommand{\pauliy}{
    \left( \begin{array}{cc}
        0 & -i \\
        i & 0
    \end{array}
    \right)
}
\newcommand{\pauliz}{
    \left( \begin{array}{cc}
        1 & 0 \\
        0 & -1
    \end{array}
    \right)
}

\newcommand{\twopartdef}[4]
{
    \left\{
        \begin{array}{ll}
            #1 & \mbox{if } #2 \\
            #3 & \mbox{if } #4
        \end{array}
        \right.
}

% Bold vectors
\newcommand{\bv}[1]{\boldsymbol{#1}}

% Column vectors
\newcommand{\colvec}[1]{\begin{bmatrix} #1 \end{bmatrix}}

\newtheorem*{theorem*}{Tétel}
\newtheorem*{def*}{Definíció}
\newtheorem*{pld*}{Példa}
\newtheorem*{megj}{Megjegyzés}

\begin{document}
\title{Kvantummechanika gyorstalpaló}
\author{Nagy Dániel}
\date{\today}

\maketitle
\vspace*{7.5cm}
\begin{figure}[h!]
    \begin{center}
    \includegraphics[width=0.4\textwidth]{./images/elte.eps}
    \end{center}
\end{figure}

\newpage


\section{Bevezetés}
A kvantummechanika formalizmusának az alapja a $\mathcal{H}$ komplex Hilbert tér. A kvantummechanikai állapotok a $\mathcal{H}$ Hilbert tér
elemei, vektorai.

\subsection{Dirac-jelölés}
A $\mathcal{H}$ Hilbert tér elemei a $\ket{\psi}$ vektorok. A $\mathcal{H}^{*}$ duális térben a vektorokat $\bra{\phi}$
-vel jelöljük.
\\ A vektorok skalárszorzására a következő igaz:
\begin{itemize}
    \item $\braket{\phi | \psi} \in \mathbb{C}$
    \item $\braket{\phi | \psi} = \braket{\psi|\phi}^{*}$
    \item $\norm\psi ^2 = \braket{\psi|\psi} \geq 0$
\end{itemize}
További azonosságok:
\begin{itemize}
    \item $(c\ket\psi)^{*} = c^* \bra\psi$
\end{itemize}

\subsubsection{Reprezentáció}
Számolások során a Hilbert-tér általában az $\mathcal{L}^2$-tér.
Ilyenkor az absztrakt $\ket\psi$ vektorok helyett a $\psi(x)$ függvényt használjuk:
\begin{equation}
    \psi: \mathbb{R} \rightarrow \mathbb{C} \in \mathcal{L}^2 \iff \norm \psi^2 = \int\limits_{-\infty}^{\infty} |\psi(x)|^2 dx < \infty
\end{equation}

\subsection{Kvantumállapot és Born-féle értelmezés}
A kvantummechanikában minden állapothoz rendelhető egy $\ket\psi$ vektor.
Az ilyen állapotok 1-re normált állapotok kell hogy legyenek, vagyis
$\norm \psi = \braket{\psi|\psi} = 1$.
Abban az esetben ha adott egy $\ket j$ \textit{diszkrét, ortonormált} bázis, $\ket \psi$ felírható ezen a bázison:
\begin{equation}
    \ket\psi = \sum\limits_{j=0}^{n} c_j\ket j = \sum\limits_{j=0}^{n} \braket{j|\psi}\ket j
\end{equation}

ahol  $c_j = \braket{j|\psi}$ komplex együtthatók. Fontos: a Born-féle értelmezés alapján $|c_j|^2$ annak a valószínűsége, hogy a $\ket\psi$ állapottal jellemzett
rendszer éppen a $\ket j$ sajátállapotban található.
A normálási feltételből következik, hogy 
\begin{equation}
    1=\braket{\psi|\psi} = \sum\limits_{k=0}^{n} c_k^* \sum\limits_{j=0}^{n} c_j \braket{k|j} = \sum\limits_{k=0}^{n} |c_k|^2 = 1
\end{equation}
továbbá
\begin{equation}
    \sum\limits_{j=0}^{n} \ket j \bra j = \hat{1}
\end{equation}

A valószínűségi értelmezés alapján, ha egy rendszer a $\ket a$ állapotban van preparálva és arra vagyok kíváncsi, hogy mi a valószínűsége annak,
hogy $\ket b$ állapotban találom, akkor a $p = |\braket{b|a}|^2$ értéket kell kiszámolni.

Ha $\ket\phi$ és $\ket\psi$ két kvantumállapot, akkor az átmeneti amplitúdó:
\begin{equation}
    \braket{\phi|\psi} = \sum\limits_j \braket{\phi|j}\braket{j|\psi} = \sum\limits_j \phi_j^{*}\psi_j
\end{equation}
\subsection{Fiziaki mennyiségek és operátorok}
A $\mathcal{H}$ Hilbert téren értelmezhetők az $\hat{A} : \mathcal{H} \rightarrow \mathcal{H}$  típusú operátorok.
Értelmezhető az $\hat{A}$ operátor hermitikus adjungáltja a követkző képpen :
\begin{equation}
    \bra \phi \hat{A} \ket \psi = \bra \psi \hat{A}^\dag \ket \phi
\end{equation}
Itt $\hat{A}^\dag$ -ot az $\hat{A}$ hermitikus adjungáltjának nevezzük és igaz, hogy ha
\begin{equation}
    \ket w = \hat{A}\ket v \implies \bra w = \bra v \hat{A}^\dag
\end{equation}

\begin{def*}
    A $\hat{H}$ operátort hermitikusnak nevezzük ha $\hat H^\dag = \hat H$.
\end{def*}

\begin{theorem*}
    A kvantummechanikában minden fizikai állapothoz rendelhető egy hermitikus operátor.
\end{theorem*}

\begin{pld*}
    Az impulzus operátora $\hat p = -i\hbar \nabla$, az energia operátora $\hat H = \frac{\hat p^2}{2m} + \hat V$.
\end{pld*}

\begin{def*}
    Egy $\hat U$ operátort unitérnek nevezünk, ha $\hat U^{-1} = \hat U ^\dag$.
\end{def*}

\begin{pld*}
    Az időfejlődés operátora $\hat U(t) = e^{-\frac{i}{\hbar}\hat H t}$ unitér operátor.
\end{pld*}

A hermitikus operátorok sajátértékei mindig valósak. Ha egy operátor sajátvektorai ortonormált bázist alkotnak, akkor:
$$ \hat A = \sum\limits_{n} a_n \ket n \bra n \textrm{, ahol } \hat A \ket n = a_n \ket n $$
\begin{theorem*}
    Egy hermitikus operátor sajátvektorai mindig ortgonálisak.
\end{theorem*}
Illetve, ha egy $\hat A$ operátor sajátvektorai ortonormált bázist alkotnak, akkor egy tetszőleges $\ket \psi$ vektor
kifejthető ezen a bázison:
$$\hat A \ket n = a_n \ket n \textrm{, és} \braket{n|m} = \delta_{mn} \implies \ket\psi = \sum\limits_n \braket{n|\psi} \ket n 
\equiv  \sum\limits_n c_n \ket n$$ 
Legyen most $\hat A$ hermitikus, $\ket\psi$ egy állapot, úgy, hogy $\hat A \ket n = a_n\ket n$ és $\ket\psi = \sum\limits_n c_n \ket n$.
Ekkor az $\hat A \ket\psi$ vektor $b_m$ együtthatói az $\ket n$ bázison:
$$b_m = \braket{m|\hat A \psi} = \sum\limits_n \bra{m} \hat A \ket n \braket{n|\psi} = \sum\limits_n A_{mn}c_n$$
Ezért $$A_{mn} = \braket{ m |\hat A | n }$$
$|\braket{ \phi | \hat A | \psi}|^2$ az minek a valószínűsége?

\subsection{Folytonos bázis, impulzus-reprezentáció}
\subsubsection{Hely-bázis}
Eddig azzal foglalkoztunk, mi van, ha diszkrét bázisunk van, pl az $\ket n$ energia-sajátállapotok bázisa. 
A kvantummechanikában a bázis általában valamilyen mérhető mennyiséghez tartozó $\hat O$ hermitikus operátor sajátvektoraiból áll.
Vannak olyan operátorok, melyeknek a spektruma folytonos, így a bázis is folytonos lesz.
A leggyakoribb ilyen eset az $\hat x$ helyoperátorhoz tartozó $\ket x$ hely-bázis és a $\hat p$ impulzus-operátorhoz
tartozó $\ket p$ impulzus-bázis.
\\
Ekkor $\ket x$ az az állapot amelyben a részecske pontosan az $x$ helyen van és $\hat x\ket x = x\ket x$.
Hasonlóan $\ket p$ is impulzus-sajátállapot: $\hat p\ket p = p\ket p$.
Az előbbiek alapján egy tetszőleges $\ket\psi$ állpotban annak a valószínűsége, hogy a részecske az $x$ helyen található:
$|\braket{x|\psi}|^2$. A $\braket{x|\psi}$ valószínűségi amplitúdót nevezzük hullámfüggvénynek:
\begin{equation}
    \psi(x)=\braket{x|\psi}
\end{equation}
A normálási feltétel így fog kinézni:
\begin{equation}
    1 = \int\dd x |\psi(x)|^2 = \int\dd x \braket{\psi|x}\braket{x|\psi}
\end{equation} 
illetve
\begin{equation}
    \int\dd x \ket x \bra x = \hat 1
\end{equation}
Hasonlóan a diszkrét esethez, itt is definiáljuk az ortonormáltságot:
\begin{equation}
    \braket{x'|x} = \delta(x-x')
\end{equation}
ezért 
\begin{equation}
    \psi(x') = \braket{x'|\psi} = \int\dd x \braket{x'|x}\braket{x|\psi} = \int\dd x \delta (x-x')\psi(x)
\end{equation}
Az átmeneti amplitúdó a $\ket\phi$ és $\ket\psi$ kvantumállapotok között így fog alakulni:
\begin{equation}
    \braket{\phi|\psi} = \int\dd x \braket{\phi|x}\braket{x|\psi} = \int\dd x \phi^{*}(x)\psi(x)
\end{equation}

\subsubsection{Impulzus-bázis}
Impulzus-bázisban igaz, hogy $\hat p \ket p = p\ket p$. Tudjuk, hogy $\hat p = -i\hbar\partial_x$ (1D-ben).
Állítsuk elő a $\ket p$ állapotot $\ket x$ bázisban:
\begin{equation}
    \ket p = \int\dd x \ket x \braket{x|p}
\end{equation}
és legyen $\braket{x|p}=\phi_p(x)$. Ekkor egyrészt igaz, hogy
\begin{equation}
    \braket{x|\hat p|p} = \braket{x|p|p} = p\braket{x|p} = p\phi_p(x)
\end{equation}
másrészt
\begin{equation}
    \hat p \braket{x|p} = -i\hbar\frac{\partial \phi_p(x)}{\partial x} 
\end{equation}
tehát
\begin{equation}
    -i\hbar\frac{\partial \phi_p(x)}{\partial x} =  p\phi_p(x) \implies \phi_p(x) = A e^{\frac{i}{\hbar} px}
\end{equation}
A normalizálás az előbbihez hasonló módon történik:
\begin{equation}
    \braket{p'|p} = \delta(p-p')
\end{equation}
emiatt:
\begin{eqnarray}
    \delta(p-p') = \braket{p'|p} = \int\dd x \braket{p'|x}\braket{x|p} \\
    = \int\dd x \phi_{p'}^{*}(x)\phi_p(x) \\
    = \int\dd x A^2 e^{\frac{i}{\hbar} (p-p') x} \\
    = 2\pi\hbar A^2 \delta (p-p')
\end{eqnarray}
Innen pedig következik, hogy $A=1/\sqrt{2\pi\hbar}$.

Át lehet váltani a hullámfüggvényt impulzustérbe is, felhasználva, hogy $\hat x\ket x = x\ket x$ és $\braket{p|x}=\tilde\phi_x(p)$
\begin{equation}
    \ket x = \int\dd p \ket p \braket{p|x}
\end{equation}
Az előző esethez hasonlóan
\begin{equation}
    \braket{p|\hat x|x} = \braket{p|x|x} = x \braket{p|x} = x\tilde\phi_x(p)
\end{equation}
Most jön egy trükk: az $\hat x$ operátor hatása impulzustérben $\hat x = i\hbar\partial_p$, ezért
\begin{equation}
    \hat x \braket{p|x} = i\hbar \frac{\partial\tilde\phi_x(p)}{\partial p}
\end{equation}
Így 
\begin{equation}
    i\hbar \frac{\partial\tilde\phi_x(p)}{\partial p} = x\tilde\phi_x(p) \implies \tilde\phi_x(p) = B e^{-\frac{i}{\hbar}px}
\end{equation}

Az előbbiekhez hasonlóan belátható, hogy $B=1/\sqrt{2\pi\hbar}$. Azt kaptuk tehát, hogy $\tilde\phi_x(p) = \phi^{*}_p(x)$ egymás komplex konjugáltjai,
amit el is vártunk, hiszen $\braket{p|x}=\braket{x|p}^{*}$ definíció szerint.

A bázistranszformációhoz szükséges összefüggések:
\begin{eqnarray}
    \braket{x|p} = \phi_p(x) = \frac{1}{\sqrt{2\pi\hbar}}e^{\frac{i}{\hbar}px} \\ [8pt]
    \braket{p|x} = \tilde\phi_x(p) = \frac{1}{\sqrt{2\pi\hbar}}e^{-\frac{i}{\hbar}px} \\ [8pt]
    \ket p = \int\dd x \ket x \braket{x|p} = \frac{1}{\sqrt{2\pi\hbar}} \int\dd x e^{\frac{i}{\hbar}px} \ket x \\ [8pt]
    \ket x = \int\dd p \ket p \braket{p|x} = \frac{1}{\sqrt{2\pi\hbar}} \int\dd p e^{-\frac{i}{\hbar}px} \ket p \\
\end{eqnarray}


\subsection{Kommutátorok}
\begin{def*}
    Az $\hat A$ és $\hat B$ kommutátora $\commut{\hat A}{\hat B} = \hat A \hat B - \hat B \hat A$. 
\end{def*}
\begin{theorem*}
    Ha $\hat A$ és $\hat B$ kommutálnak, azaz $\commut{\hat A}{\hat B} = 0$, akkor van közös sajátvektor-rendszerük, azaz létezik olyan
    $\psi$ állapot mindkét operátornak sajátállapota.
\end{theorem*}
\begin{theorem*}
    Ha $\hat A$ és $\hat B$ \textbf{nem} kommutálnak, azaz $\commut{\hat A}{\hat B} \neq 0$, akkor \textbf{nincs} közös sajátvektor-rendszerük, azaz nem 
    létezik olyan $\ket\psi$ állapot, amely mindkét operátor sajátállapota.
\end{theorem*}
\begin{megj}
    Ha két operátor nem kommutál, attól még egy adott $\ket\psi$ állapot mindkét operátor sajátvektorainak bázisán kifejthető lehet,
    és ezek egymásba áttranszformálhatók. Pl. $\commut{\hat x}{\hat p} = i\hbar$, de létezik $\braket{p|\psi}$ és $\braket{x|\psi}$
    kifejtés is.
\end{megj}
\begin{itemize}
    \item $\commut{\hat x_i}{\hat x_j} = 0$
    \item $\commut{\hat p_i}{\hat p_j} = 0$
    \item $\commut{\hat x_i}{\hat p_j} = i\hbar \delta_{ij}\mathbf 1$
    \item $\commut{\hat L_i}{\hat L_j} = i \hbar \varepsilon_{ijk}\hat L_k$
    \item $\commut{\hat S_i}{\hat S_j} = i \hbar \varepsilon_{ijk}\hat S_k$
    \item $\commut{\hat J_i}{\hat J_j} = i \hbar \varepsilon_{ijk}\hat J_k$
    \item $\commut{\hat J^2}{\hat J_i} = \commut{\hat S^2}{\hat S_i} = \commut{\hat L^2}{\hat L_i} = 0$
    \item $\commut{\hat L_i}{\hat S_j} = 0$
    \item $\commut{\hat x_i}{\hat L_j} = i \hbar \varepsilon_{ijk}\hat x_k$
    \item $\commut{\hat p_i}{\hat L_j} = i \hbar \varepsilon_{ijk}\hat p_k$    
    \item Ha $\hat A$ hermitikus $\implies \hat A^\dag = \hat A \implies \commut{\hat A^\dag}{\hat A} = 0$  
    
    
\end{itemize}
\section{Harmonikus oszcillátor}
A harmonikus oszcillátor esetén a Hamilton-operátor ugyanz, mint a klasszikus Hamilton-függvény:
\begin{equation}
  \hat H = \frac{\hat p^2}{2m} + \frac{1}{2}m\omega^2 \hat x^2  
\end{equation}
3D esetben:
\begin{equation}
    \hat H = \frac{1}{2m}(\hat p_x^2+\hat p_y^2+\hat p_z^2) + \frac{1}{2}m\omega^2 (\hat x^2 + \hat y^2 + \hat z^2)  
\end{equation}
Az energiasajátértékekhez meg kell oldani a $\hat H \ket n = E_n \ket n$ egyenletet. Ennek a megoldása 1D-ben
folytonos bázison:
\begin{equation}
    \ket n = \psi_n(x) = \frac{1}{\sqrt{(2^n n!)}} \cdot \left (\frac{m\omega}{\pi\hbar}\right)^{1/4}
    \cdot e^{-\frac{m\omega x^2}{2\hbar}} \cdot H_n\left(\sqrt\frac{m\omega}{\hbar}x\right)    
\end{equation}
ahol $H_n$-ek a Hermite-polinomok:
\begin{equation}
    H_n(z) = (-1)^n e^z^2 \frac{\dd^n}{\dd z^n}e^{-z^2}
\end{equation}
Ezekből 
\begin{equation}
    E_n = \left(n + \frac{1}{2}\right)   \hbar \omega 
\end{equation}
\pagebreak
Az állapotok léptető operátorait a következő képpen definiáljuk:
\begin{itemize}
    \item $
        \hat a = \sqrt{\frac{m\omega}{2\hbar}}\left(\hat x + \frac{i}{m\omega}\hat p\right)
    $
    \item $
       \hat a^{\dag} = \sqrt{\frac{m\omega}{2\hbar}}\left(\hat x - \frac{i}{m\omega}\hat p\right)
    $
    \item $
        \commut{\hat a}{\hat a^{\dag}} = \hat 1
    $
    \item $
        \hat a \ket n = \sqrt n \ket{n-1}
    $
    \item $
        \hat a^{\dag} \ket n = \sqrt{n+1}\ket{n+1}
    $
    \item $
        \ket n = \frac{(\hat a^{\dag})^n}{\sqrt{n!}}\ket 0 
        \textrm{  folytonos bázison  } \ket 0 = Ce^{-\frac{m\omega x^2}{2\hbar}}
    $
\end{itemize}


\section{Impulzusmomentumok}
\subsection{Pálya-impulzusmomentum}
A kvantummechanikai impulzusmomentum operátora $\hat{\vec L} = \hat{\vec r} \times \hat{\vec p} = -i \hbar (\vec r \times \nabla)$.

Gömbi koordinátarendszerben felírva:
\begin{itemize}
    \item $
        \hat L^2 = -\hbar^2 \left(\frac{1}{\sin \theta} \frac{\partial}{\partial \theta} \left(\sin\theta\frac{\partial}{\partial \theta}\right) 
        + \frac{1}{\sin^2 \theta}\frac{\partial^2}{\partial\varphi^2}\right)
    $
    \item $
        \hat L_z = -i\hbar \frac{\partial}{\partial\varphi}
    $
    \item $
      \hat L_{\pm} = \hat L_x \pm i \hat L_y
    $
    \item $
        \hat L_+ = \hbar e^{i\varphi} \left( \frac{\partial}{\partial \theta} - i \cot\theta \frac{\partial}{\partial \varphi}\right)
    $
    \item $
        \hat L_- = \hbar e^{-i\varphi} \left( - \frac{\partial}{\partial \theta} - i \cot\theta \frac{\partial}{\partial \varphi}\right)
    $
    \item $
        \commut{\hat L_+}{\hat L_-} = 2 \hbar L_z
    $
    \item $
        \hat L^2 = \hat L_+ \hat L_- + \hat L_z^2 - \hbar \hat L_z
    $
\end{itemize}
Az $\hat L$ operátor sajátállapotait két kvantumszám határozza meg , $l$ és $m_l$, a sajátvektorokat ezért $\ket{l,m_l}$ jelöli.
A sajátértékek a következők lehetnek: $l\in \{0,1,2, ...\}$, $m_l \in \{-l,... 0, ..., l\}$.

Ezzel a jelöléssel igazak a következő azonosságok:
\begin{itemize}
    \item $
        \hat L^2 \ket{l, m_l} = \hbar^2 l (l+1) \ket{l,m_l}
    $
    \item $
        \hat L_z \ket{l, m_l} = \hbar m_l \ket{l,m_l}
    $
    \item $
      \hat L_{\pm}\ket{l, m_l} = \hbar \sqrt{l(l+1) - m_l(m_l\pm1)}\ket{l,m_l\pm 1}
    $
    \item $
        \hat L_+ \ket{l,m_l=l} = 0 ~~ \textrm{és} ~~ \hat L_- \ket{l,m_l=-l} = 0
    $
\end{itemize}
Ha $\ket{l,m_l}$-eket az $\mathcal{L}^2$ téren akarjuk ábrázolni, akkor:
\begin{equation}
    \braket{\theta, \varphi | l,m_l} = Y^{m_l}_l(\theta, \varphi)
\end{equation}
\begin{equation}
   \braket{l',m'_l|l,m_l} = \int \limits_0^{\pi}\sin\theta d \theta \int \limits_0^{2\pi}d \varphi (Y^{m_l'}_{l'})^{*} Y^{m_l}_l = \delta_{ll'}\delta_{m_lm_l'}
\end{equation}
Az $Y^{m_l}_l(\theta, \varphi)$ függvények a \ndurl{https://en.wikipedia.org/wiki/Spherical_harmonics}{gömbi harmonikusok}

\subsection{Spin}
A spin-operátort $\hat S$ jelöli, sajátállapotait két kvantumszám adja meg, $s$ és $m_s$, úgy, \\ hogy
$s\in\{0, \frac{1}{2}, 1, \frac{3}{2}, ...\}$ és $m_s \in \{-s, -s+1,...,s-1, s\}$

\subsubsection{Az $s = \frac{1}{2}$ spin és a Pauli-mátrixok}
Feles spin esetén $s=\frac{1}{2}$, ezért $m_s \in \{-\frac{1}{2},\frac{1}{2}\}$. Az állapotokat a következő rövidítéssel jelöljük:
\begin{equation}
    \ket + \equiv \ket{\frac{1}{2}, \frac{1}{2}}    
\end{equation}
\begin{equation}
    \ket - \equiv \ket{\frac{1}{2}, -\frac{1}{2}}    
\end{equation}
Ez egy kétállapotú rendszert határoz meg, egy tetszőleges állapotvektor:
\begin{equation}
    \ket\psi = \alpha\ket+ + \beta\ket- ~,~~ |\alpha|^2 + |\beta|^2 = 1    
\end{equation}
Ebben az esetben a spin operátora $\vec S = \frac{\hbar}{2}\vec{\sigma}$, ahol
\begin{equation} \sigx = \paulix \end{equation}
\begin{equation} \sigy = \pauliy \end{equation}
\begin{equation} \sigz = \pauliz \end{equation}

A $\sigma_k$-ra vonatkozó azonosságok:
\begin{itemize}
    \item $\sigz \ket+ = \frac{\hbar}{2}\ket+$
    \item $\sigz \ket- = -\frac{\hbar}{2}\ket-$    
    \item $
        \commut{\sigma_i}{\sigma_j} = 2i\varepsilon_{ijk}\sigma_k
    $
    \item $
        \sigma_k\sigma_l = \delta_{kl}\mathbf 1 + i\varepsilon_{klj}\sigma_j
    $
    \item $
        (\vec u \vec \sigma)(\vec v\vec\sigma) = (\vec u\cdot\vec v)\mathbf 1 + i(\vec u \times\vec v)\vec\sigma
    $
    \item $
        \sigma_+ = \sigx + i\sigy
    $
    \item $
        \sigma_- = \sigx - i\sigy
    $

\end{itemize}
\subsection{Teljes impulzusmomentum}
\subsection{Impulzusmomentumok összeadása és Clebsch–Gordan együtthatók}
Az pálya-impulzusmomentum és a spin különböző Hilbert-tér elemei, ezért ha össze akarom adni őket,
akkor a Hilbert-terek direkt szorzatát kell képezni.
Legyen egy rendszer (pl. H-atom elektronja) ami spinnel és pálya-impulzusmomentummal is rendelkezik.
Legyenek a spin sajátállapotok $\ket{s, m_s}$ illetve a pálya-impulzusmomentum sajátállapotok $\ket{l, m_l}$.
Ekkor a teljes impulzusmomentum sajátállapotai a két tér elemeinek direkt szorzata:
\begin{equation}
    \ket{j, m} = \ket{l, m_l} \otimes \ket{s, m_s} \equiv \ket{l, m_l}\ket{s, m_s} \equiv \ket{l, m_l,s, m_s}
\end{equation} 

Ezen a szorzattéren hat a $\hat{\vec J}$ operátor, amit így definiálunk:
\begin{equation}
    \hat{\vec J} = \hat{\vec L} + \hat{\vec S}
\end{equation}
Ez a jelölés pongyola, mert a $\hat{\vec J}$-nek a szorzattéren vagyis $\mathcal{H}_L \otimes \mathcal{H}_S$-en kell hatnia,
ezért korrektül:
\begin{equation}
    \hat J_k = \hat L_k \otimes \hat 1 + \hat 1 \otimes \hat S_k
\end{equation}
Ez a következő képpen hat a sajátállapotokra:
\begin{equation}
    \hat J^2 \ket{j, m} = \hbar^2 j(j+1)\ket{j,m}
\end{equation}
\begin{equation}
    \hat J_z \ket{j, m} = \hbar m\ket{j,m}
\end{equation}
Szintén be lehet vezetni a léptető operátorokat:
\begin{itemize}
    \item $
        \hat J_{\pm} = \hat J_x \pm  i \hat J_y
    $
    \item $
       \hat J_+ \hat J_-  = \hat J ^2 - \hat J_z^2 + \hbar \hat J_z
    $
    \item $
       \hat J_- \hat J_+  = \hat J ^2 - \hat J_z^2 - \hbar \hat J_z
    $
    \item $
        \commut{\hat J_+}{\hat J_-} = 2 \hbar \hat J_z
    $
    \item $
        \commut{\hat J_z}{\hat J_{\pm}} = \pm \hbar \hat J_{\pm}
    $
    \item $
        \hat J_{\pm}\ket{j,m} = \hbar \sqrt{j(j+1) - m(m\pm1)}\ket{j,m\pm1}
    $
\end{itemize}
\section{Hidrogén-atom}
\begin{megj}
    Ha a főkvantumszám $n$, akkor $l_{max}=n-1$ vagyis $l \in \{0,1,...,n-1\}$
\end{megj}
\section{Perturbációszámítás}
\subsection{Időfüggetlen perturbációszámítás}
\subsubsection{Nemdegenerált eset}
Adott egy $\hat H$ operátor ami úgy néz ki, hogy $\hat H = \hat H^{(0)} + \lambda \hat W$, ahol $\lambda$ kicsi, 
$\lambda\hat W$ a perturbáció, $\hat H^{(0)}$-nak pedig ismerjük a sajátállapotait:
\begin{equation}
    \hat H^{(0)} \ket{n^{(0)}} = E_n^{(0)} \ket{n^{(0)}} ~, ~~ n=1,2,... 
\end{equation}
A $\hat H$ sajátérték-problémájának megoldása így fog kinézni sorfejtett alakban:
\begin{equation}
    E_n = \lambda^0 E_n^{(0)} + \lambda^1 E_n^{(1)} + \lambda^2 E_n^{(2)} + ...
\end{equation}
\begin{equation}
  \ket n = \lambda^0 \ket{n^{(0)}} + \lambda^1 \ket{n^{(1)}} + \lambda^2 \ket{n^{(2)}} + ...  
\end{equation}
Az elsőrendű korrekciók a fentiek alapján:
\begin{equation}
    E_n^{(1)} = \bra{n^{(0)}} \hat W \ket{n^{(0)}}      
\end{equation}
\begin{equation}
    \ket{n^{(1)}} = \sum\limits_{k\neq n}\frac{\bra{k^{(0)}} \hat W \ket{n^{(0)}}}{E_n^{(0)}-E_k^{(0)}}
    \ket{k^{(0)}}
\end{equation}

Másodrendű energia-korrekció:
\begin{equation}
    E_n^{(2)} = \sum\limits_{k\neq n}\frac{\left|\bra{k^{(0)}} \hat W \ket{n^{(0)}}\right|^2}{E_n^{(0)}-E_k^{(0)}}    
\end{equation}

\textbf{Figyelem!} A fenti korrekciók csak nemdegenerált energiaspektrum esetében igazak.

\subsubsection{Degenerált eset}
Akkor kell alkalmazni, amikor két vagy több energiaszint egybeesik, vagyis létezik egy degenerált sajátaltér.
Ekkor a következő szerint járunk el:
\begin{itemize}
    \item Vegyük a degenerált altérbe tartozó \textbf{perturbálatlan sajátvektorokat}. 
    \item Ezután ezekkel a vektorokkal számítsuk ki a $W_{mn} = \bra m \hat W \ket n$ mátrixelemeket.
    \item Az így kiszámolt mátrix sajátérték-problémájának megoldásai lesznek a perturbált energiaszintek.
\end{itemize}
\subsection{Időfüggő perturbációszámítás}
Most megengedjük, hogy a perturbáció függjön expliciten az időtől:
\begin{equation}
    \hat H(t) = \hat H^{(0)} + \lambda\hat W(t)
\end{equation}
A következő lépésben be kell vezetni az $\omega_{mn} = (E_m-E_n)/\hbar$ mennyiséget, és a 
perturbálatlan $\hat H^{(0)}$ operátor $\ket m$ illetve $\ket n$ sajátvektoraiból kiszámítjuk
a $W_{mn}(t) = \bra m \hat W(t)\ket n$ mátrixelemeket. Legyen most $\ket i$ a kezdeti (initial) $\ket f$
pedig a végső (final) állapot. Ekkor annak a valószínűsége, hogy $\ket i$ állapotból indulva $t$
idő elteltével a rendszer $\ket f$ állapotban található:
\begin{equation}
    P(i\rightarrow f) = P_{fi} = \frac{1}{\hbar^2}\left|\int\limits_0^t W_{fi}(\tau) e^{i\omega_{fi}\tau} d\tau\right|^2   
\end{equation}
\end{document}
